\documentclass[a4paper, 12pt]{book}
% preamble.tex

\usepackage[utf8]{inputenc}
\usepackage[T1]{fontenc}
\usepackage[english]{babel} % Ho cambiato in italian, supponendo che la maggior parte del testo sia in italiano
\usepackage{amsmath}
\usepackage{amssymb}
\usepackage{graphicx}
\usepackage{hyperref}
\usepackage{geometry}
\usepackage{fontspec} % Utile per Xetex/LuaLaTeX ma mantenuto
\usepackage{booktabs}
\usepackage{amsthm}
\usepackage{enumitem}
\usepackage{caption}
\usepackage{thmtools}
\usepackage{standalone} % Già presente nel tuo elenco
\usepackage{tkz-euclide}
\usepackage{graphicx}
\usepackage[T1]{fontenc}
\usepackage{wesa}

% --- PACCHETTI AGGIUNTI PER LA FORMATTAZIONE ---
\usepackage{parskip}  % Aggiunge spazio tra i paragrafi (stile "blocco")
\usepackage{fancyhdr} % Per intestazioni e piè di pagina personalizzati

% --- IMPOSTAZIONI GEOMETRIA E FONT ---
\geometry{a4paper, margin=1in}

% Configura lo stile 'plain' (usato da \maketitle e \tableofcontents)
% per essere coerente con il piè di pagina
\fancypagestyle{plain}{
    \fancyhf{} % Pulisce
    \cfoot{\thepage} % Mette solo il numero di pagina
    \renewcommand{\headrulewidth}{0pt} % Niente linea su queste pagine
}

% Configurazione Header/Footer per le altre pagine
\pagestyle{fancy}
\fancyhead{} % Pulisce l'intestazione
\fancyfoot{} % Pulisce il piè di pagina
\fancyhead[L]{\nouppercase{\rightmark}} % Nome Sezione a sinistra (o Chapter)
\fancyhead[R]{\nouppercase{\leftmark}}  % Nome Capitolo a destra (o Sezione)
\fancyfoot[C]{\thepage} % Numero di pagina al centro
\renewcommand{\headrulewidth}{0.4pt}
\renewcommand{\footrulewidth}{0pt}


% --- DEFINIZIONI TEOREMI ---
% (Questa parte è cruciale per la tua struttura)
\newtheoremstyle{mythmstyle}%
    {}%
    {}%
    {\itshape}%
    {}%
    {\bfseries}%
    {}%
    { }%
    {\thmname{#1}\thmnumber{ #2}%
    \thmnote{: #3\addcontentsline{toc}{subsubsection}{{\itshape#1}: #3}}. }

% 2. Imposta questo stile come quello ATTUALE
\theoremstyle{mythmstyle}

\newtheorem{theorem}{Th}[section]
\newtheorem{definition}{Def}[section]
\newtheorem{proposition}{Prop}[section]
\newtheorem{corollary}{Cor}[section]
\newtheorem{algorithm}{Alg}[section]
\newtheorem{lemma}{Lemma}[section]
\newtheorem*{remark}{Rem}

% Definizioni utili per il Main file
\newcommand{\inputlesson}[1]{\clearpage\input{#1}\clearpage} % Comando per input lezione
\usepackage[x11names]{xcolor}
\usepackage{tikz}
\usetikzlibrary{shapes,fit}
\usepackage{amsmath} % Required for align environment


\begin{document}
\setcounter{chapter}{6}
\chapter{Lesson 28-11-2025}
\section{Orthogonal Frequency Division Multiplexing (OFDM)}

\subsection{OFDM Transmitter}

\subsubsection{Notation Change and "Blockization"}
In OFDM, we change the standard notation compared to single-carrier systems. Usually, input symbols are indicated as $x[n]$ in the time domain.
In OFDM, we denote the input symbols as $x[k]$. The $k$ represents the index in the \textbf{frequency domain}, since these symbols will be transmitted on different subcarriers.

The input symbols are grouped ("packed") into blocks of size $K$.
\begin{align}
\mathbf{x} = [x[0], x[1], \dots, x[K-1]]
\end{align}
These symbols may belong to a complex constellation, for example, 256-QAM.

\noindent\colorbox{yellow!25}{
    \parbox{\dimexpr\linewidth-2\fboxsep}{
        \textbf{\faQuestionCircle\ Student Q\&A: Block Stretching} \\
        \small
        \textbf{Q:} What is meant by "stretching" the blocks?
    }
}

\vspace{0.5em} % Add a small vertical space between the question and answer

\noindent\colorbox{green!15}{
    \parbox{\dimexpr\linewidth-2\fboxsep}{
        \textbf{\faCheck\ Answer: Block Stretching} \\
        \small
        \textbf{A:} During the lecture, a student asked for clarification on "lengthening" the blocks. The professor explained that to transmit a block of $K$ symbols, if one wants to maintain the same input rate, time or bandwidth must be increased. However, in the OFDM perspective, the key concept is that we take $K$ symbols and transmit them \textit{in parallel} over a longer time period compared to a single symbol, occupying a bandwidth $B \approx K/T$ (where $T$ is the block time). Technically, in the context of the explanation, "stretch" referred to the insertion of the Cyclic Prefix, which physically lengthens the duration of the transmitted block from $K$ to $K+L$ samples.
    }
}

\subsubsection{IDFT (Inverse Discrete Fourier Transform)}
The fundamental operation at the transmitter is the IDFT applied to each block of $K$ symbols. This step brings the symbols from the frequency domain to the time domain.
The formula to generate the signal in the time domain $x[n]$ is:
\begin{align}
x[n] = \frac{1}{K} \sum_{k=0}^{K-1} x[k] e^{j \frac{2\pi}{K} nk} \quad \text{for } n=0, \dots, K-1
\end{align}

\subsubsection{Cyclic Prefix (CP) Insertion}
To mitigate multipath channel effects, a Cyclic Prefix is inserted. If the length of the channel impulse response is $L$, we insert a prefix of length $L$ (or greater) at the beginning of the block.
The extended vector $\overline{x}$ will have dimension $K+L$. The operation consists of copying the last $L$ samples of the block $x[n]$ and placing them at the beginning:
\begin{align}
\overline{x} = [\underbrace{x[K-L], \dots, x[K-1]}_{\text{Cyclic Prefix (copy)}}, \underbrace{x[0], x[1], \dots, x[K-1]}_{\text{Original Block}}]
\end{align}
Formally, the negative indices in the prefix correspond to the final indices of the original block: $x[-1] \leftarrow x[K-1]$, $x[-2] \leftarrow x[K-2]$, and so on.

\textbf{Cyclic Property:}
The insertion of the CP makes the signal "circular" with respect to the block length $K$. Mathematically, thanks to the periodicity of the complex exponential, the following relationship holds:
\begin{align}
x[-1] = x[K-1]
\end{align}
In fact, $e^{j\frac{2\pi}{K}(-1)k} = e^{j\frac{2\pi}{K}(K-1)k}$.

\section{The Channel}

The transmitted signal passes through a channel modeled as a Finite Impulse Response (FIR) filter with coefficients $h[l]$ of length $L+1$ ($l=0 \dots L$).
The channel output $y[n]$ is given by the linear convolution between the transmitted signal (including the CP) and the channel response:
\begin{align}
y[n] = \sqrt{G} \sum_{l=0}^{L} h[l] \overline{x}[n-l]
\end{align}

\textbf{Inter-Block Interference (IBI):}
For indices corresponding to the cyclic prefix (e.g., $n = -1, \dots, -L$), the convolution "picks up" symbols from the previous transmitted block. For example, for $y[-1]$, the summation involves $\overline{x}[-1-L]$, which belongs to the previous block.

\noindent\colorbox{yellow!25}{
    \parbox{\dimexpr\linewidth-2\fboxsep}{
        \textbf{\faQuestionCircle\ Student Q\&A: CP Constraints} \\
        \small
        Is there a constraint on the length of the Cyclic Prefix? \\
    }
}

\vspace{0.5em} % Add a small vertical space between the question and answer

\noindent\colorbox{green!15}{
    \parbox{\dimexpr\linewidth-2\fboxsep}{
        \textbf{\faCheck\ Answer: CP Constraints} \\
        \small
        Yes. The length of the Cyclic Prefix ($L$) must be greater than or equal to the length of the channel impulse response. This ensures that all Inter-Block Interference (IBI) falls within the cyclic prefix, which is then discarded at the receiver, leaving the data block "clean".
    }
}

\section{The OFDM Receiver}

The receiver must recover the original symbols $x[k]$. The operations are specular to those of the transmitter.

\subsection{Cyclic Prefix Removal}
The first $L$ received samples ($y[-L], \dots, y[-1]$) are contaminated by interference from the previous block. The solution is simple: discard the cyclic prefix.
We consider only the vector $y$ of length $K$:
\begin{align}
y = [y[0], y[1], \dots, y[K-1]]
\end{align}

\subsection{DFT (Discrete Fourier Transform)}
On the vector cleaned of the cyclic prefix, we apply the DFT:
\begin{align}
\tilde{y}[n] = \sum_{p=0}^{K-1} y[p] e^{-j \frac{2\pi}{K} np}
\end{align}
Here $n$ represents the subcarrier index at reception.

\section{Mathematical Proof [The "Heart" of OFDM]}

Let us see why OFDM works by substituting the channel expression into the DFT formula.

\subsection{Step 1: Expansion of $\tilde{y}[n]$}
Substitute $y[p]$ with the channel convolution:
\begin{align}
\tilde{y}[n] = \sum_{p=0}^{K-1} \left( \sqrt{G} \sum_{l=0}^{L} h[l] \overline{x}[p-l] \right) e^{-j \frac{2\pi}{K} np}
\end{align}

\subsection{Step 2: Substitution of $x$ (IDFT)}
Substitute $\overline{x}[p-l]$ with its IDFT definition (the original input $x[k]$):
\begin{align}
\tilde{y}[n] = \sum_{p=0}^{K-1} \sqrt{G} \sum_{l=0}^{L} h[l] \left( \frac{1}{K} \sum_{k=0}^{K-1} x[k] e^{j \frac{2\pi}{K} (p-l)k} \right) e^{-j \frac{2\pi}{K} np}
\end{align}
\textit{Note:} Thanks to the cyclic prefix, the time shift $p-l$ in the linear convolution acts as a circular shift, allowing the use of the IDFT complex exponentials.

\subsection{Step 3: Rearrangement of Summations}
We bring out the constants and group the terms:
\begin{align}
\tilde{y}[n] = \frac{\sqrt{G}}{K} \sum_{k=0}^{K-1} x[k] \left( \sum_{l=0}^{L} h[l] e^{-j \frac{2\pi}{K} lk} \right) \left( \sum_{p=0}^{K-1} e^{j \frac{2\pi}{K} p(k-n)} \right)
\end{align}

Let us analyze the two terms in parentheses:
\begin{enumerate}
    \item \textbf{Channel DFT ($\tilde{h}[k]$):} The term $\sum_{l=0}^{L} h[l] e^{-j \frac{2\pi}{K} lk}$ is exactly the frequency response of the channel at frequency $k$, which we call $\tilde{h}[k]$. We assume $h[l]=0$ for $l>L$.
    \item \textbf{Orthogonality:} The term $\sum_{p=0}^{K-1} e^{j \frac{2\pi}{K} p(k-n)}$ sums complex exponentials. This sum is non-zero only if $k=n$.
    \begin{align}
    \sum_{p=0}^{K-1} e^{j \frac{2\pi}{K} p(k-n)} = K \cdot \delta(k-n)
    \end{align}
    where $\delta(x) = 1$ if $x=0$, otherwise $0$.
\end{enumerate}

\subsection{Final Result}
Thanks to the Dirac delta $\delta(k-n)$, the summation over $k$ "survives" only for the term where $k=n$. The factor $K$ in the numerator cancels with $1/K$. We obtain:
\begin{align}
\tilde{y}[n] = \sqrt{G} \cdot \tilde{h}[n] \cdot x[n]
\end{align}

\noindent\colorbox{yellow!25}{
    \parbox{\dimexpr\linewidth-2\fboxsep}{
        \textbf{\faQuestionCircle\ Student Q\&A: Notation Confusion} \\
        \small
        Confusion regarding $n$ vs $k$ notation in the final proof. \\
    }
}

\vspace{0.5em}
\noindent\colorbox{green!15}{
    \parbox{\dimexpr\linewidth-2\fboxsep}{
        \textbf{\faCheck\ Answer: Notation Confusion} \\
        \small
        Some ambiguity was noted in the use of indices. In the final formula, $n$ is used as the output subcarrier index ($\tilde{y}(n)$). The formula proves that the output at index $n$ depends \textit{only} on the input at the same index $n$ ($x[n]$), scaled by the channel $h[n]$. This confirms the absence of Inter-Carrier Interference (ICI).
    }
}

\subsection{Interpretation and Equalization}
The result $\tilde{y}[n] = \sqrt{G} \tilde{h}[n] x[n]$ is fundamental. It means that, in the frequency domain, there is no interference between symbols or subcarriers. The symbol transmitted on frequency $n$ ($x[n]$) is merely scaled by a complex coefficient $\tilde{h}[n]$ (the channel response at that frequency).

\textbf{One-Tap Equalization:}
To recover the transmitted symbol $\hat{x}[n]$, a simple division (Zero-Forcing) is sufficient, assuming the channel is known:
\begin{align}
\hat{x}[n] = \frac{\tilde{y}[n]}{\sqrt{G} \cdot \tilde{h}[n]}
\end{align}
\end{document}